\section{Modificadores: Parte 1}

Los \textbf{modificadores} son cláusulas y operadores que se
combinan con \texttt{SELECT} para refinar, filtrar y ordenar los
resultados de una consulta.
En esta sección cubrimos los seis modificadores fundamentales que
todo programador SQL debe dominar desde el primer día.
Todos los ejemplos siguen usando la tabla \texttt{empleados}
introducida en la sección anterior.

%--------------------------------------------------------------
\subsection{\texttt{DISTINCT}}
%--------------------------------------------------------------

Por defecto, \texttt{SELECT} devuelve \textbf{todas} las filas que
satisfacen la consulta, incluyendo duplicados.
La palabra clave \texttt{DISTINCT} instruye al SGBD a eliminar las
filas repetidas del resultado, conservando sólo valores únicos.

\begin{definicion}{\texttt{DISTINCT}}
Modificador que se coloca inmediatamente después de \texttt{SELECT}
y elimina del resultado todas las filas cuyo conjunto de valores en
las columnas proyectadas sea idéntico a otra fila ya presente.
\end{definicion}

\begin{lstlisting}[caption={DISTINCT sobre una columna.}]
-- ¿Qué departamentos existen en la empresa?
-- Sin DISTINCT aparecerían 10 filas (una por empleado)
SELECT DISTINCT departamento
FROM   empleados;

-- Resultado:
-- Ventas
-- TI
-- RRHH
-- Finanzas
\end{lstlisting}

\begin{lstlisting}[caption={DISTINCT sobre múltiples columnas.}]
-- Combinaciones únicas de departamento + cargo
SELECT DISTINCT departamento, cargo
FROM   empleados
ORDER BY departamento, cargo;
\end{lstlisting}

\begin{nota}
\texttt{DISTINCT} opera sobre la \textit{combinación completa} de
todas las columnas proyectadas, no columna por columna.
En el ejemplo anterior, \texttt{('TI', 'Desarrolladora')} es una
combinación única aunque haya dos desarrolladoras en TI, porque
ambas tienen exactamente los mismos valores en esas dos columnas
— en ese caso sí se eliminará el duplicado.
\end{nota}

%--------------------------------------------------------------
\subsection{\texttt{WHERE}}
%--------------------------------------------------------------

La cláusula \texttt{WHERE} es el mecanismo de \textbf{filtrado por
fila} de SQL.
Evalúa una condición booleana para cada fila de la tabla fuente;
sólo las filas donde la condición resulta \texttt{TRUE} pasan al
resultado.

\begin{definicion}{\texttt{WHERE}}
Cláusula que restringe las filas devueltas por una consulta a
aquellas que satisfacen una expresión booleana dada.
Se evalúa \textit{antes} de \texttt{GROUP BY}, \texttt{HAVING} y
\texttt{SELECT}.
\end{definicion}

\begin{lstlisting}[caption={Filtros básicos con WHERE.}]
-- Empleados del departamento de TI
SELECT nombre, cargo, salario
FROM   empleados
WHERE  departamento = 'TI';

-- Empleados con salario mayor a 45,000
SELECT nombre, salario
FROM   empleados
WHERE  salario > 45000;

-- Empleados que ingresaron después del 1 de enero de 2020
SELECT nombre, fecha_ingreso
FROM   empleados
WHERE  fecha_ingreso > '2020-01-01';
\end{lstlisting}

\begin{ejemplo}
Imagina la tabla como un tamiz: cada fila cae por los agujeros a
menos que cumpla la condición de \texttt{WHERE}.
Las que no la cumplen simplemente desaparecen del resultado.
El dato original en la tabla \textbf{no se modifica}; sólo se
filtra lo que se muestra.
\end{ejemplo}

%--------------------------------------------------------------
\subsection{\texttt{ORDER BY}}
%--------------------------------------------------------------

\texttt{ORDER BY} define el \textbf{criterio de ordenamiento} del
resultado final.
Sin él, el SGBD no garantiza ningún orden particular entre
ejecuciones.

\begin{lstlisting}[caption={Orden ascendente y descendente.}]
-- Orden ascendente (A → Z, menor → mayor): por defecto
SELECT nombre, salario
FROM   empleados
ORDER BY salario ASC;   -- ASC es opcional, es el valor por defecto

-- Orden descendente (mayor → menor)
SELECT nombre, salario
FROM   empleados
ORDER BY salario DESC;
\end{lstlisting}

\begin{lstlisting}[caption={Ordenamiento compuesto por múltiples columnas.}]
-- Primero por departamento (A→Z),
-- dentro de cada departamento por salario (mayor→menor)
SELECT nombre, departamento, salario
FROM   empleados
ORDER BY departamento ASC,
         salario      DESC;
\end{lstlisting}

\begin{lstlisting}[caption={ORDER BY usando posición numérica de columna.}]
-- Se puede referenciar la columna por su posición en SELECT
-- (poco recomendado: hace el código menos legible)
SELECT nombre, departamento, salario
FROM   empleados
ORDER BY 3 DESC;   -- 3 = tercera columna = salario
\end{lstlisting}

\begin{nota}
El uso de posición numérica en \texttt{ORDER BY} está permitido
pero se desaconseja en código de producción: si el orden de las
columnas en \texttt{SELECT} cambia, el ordenamiento cambia sin
ningún aviso de error.
\end{nota}

%--------------------------------------------------------------
\subsection{\texttt{LIKE}}
%--------------------------------------------------------------

\texttt{LIKE} permite comparar una columna de texto contra un
\textbf{patrón}, usando caracteres comodín.
Es la herramienta de búsqueda de texto parcial más básica de SQL.

\begin{table}[H]
\centering
\small
\renewcommand{\arraystretch}{1.4}
\begin{tabularx}{\linewidth}{@{} c l X @{}}
\toprule
\textbf{Comodín} & \textbf{Significado} & \textbf{Ejemplo} \\
\midrule
\texttt{\%} & Cero o más caracteres cualesquiera
            & \texttt{'Ana\%'} coincide con \texttt{Ana}, \texttt{Analista}, \texttt{Anabel} \\
\texttt{\_} & Exactamente un carácter cualquiera
            & \texttt{'\_na'} coincide con \texttt{Ana}, \texttt{Ina}, \texttt{Una} \\
\bottomrule
\end{tabularx}
\caption{Comodines del operador \texttt{LIKE}.}
\label{tab:comodines-like}
\end{table}

\begin{lstlisting}[caption={Patrones con LIKE.}]
-- Nombres que empiezan con 'A'
SELECT nombre FROM empleados
WHERE  nombre LIKE 'A%';

-- Nombres que terminan con 'ez'
SELECT nombre FROM empleados
WHERE  nombre LIKE '%ez';

-- Nombres que contienen 'ar' en cualquier posición
SELECT nombre FROM empleados
WHERE  nombre LIKE '%ar%';

-- Nombres cuya segunda letra es 'o'
SELECT nombre FROM empleados
WHERE  nombre LIKE '_o%';

-- Nombres que NO contienen 'a' (NOT LIKE)
SELECT nombre FROM empleados
WHERE  nombre NOT LIKE '%a%';
\end{lstlisting}

\begin{nota}
En MySQL, \texttt{LIKE} es \textbf{insensible a mayúsculas/minúsculas}
por defecto cuando la columna usa un \textit{collation} insensible
(\texttt{utf8\_general\_ci}, \texttt{utf8mb4\_unicode\_ci}, etc.).
Si necesitas distinguir mayúsculas usa \texttt{LIKE BINARY} o cambia
el collation de la columna.
\end{nota}

%--------------------------------------------------------------
\subsection{\texttt{AND}, \texttt{OR}, \texttt{NOT}}
%--------------------------------------------------------------

Los operadores lógicos permiten \textbf{combinar múltiples
condiciones} dentro de una cláusula \texttt{WHERE}.

\begin{table}[H]
\centering
\small
\renewcommand{\arraystretch}{1.5}
\begin{tabularx}{\linewidth}{@{} l X X @{}}
\toprule
\textbf{Operador} & \textbf{Resultado TRUE cuando\ldots} & \textbf{Precedencia} \\
\midrule
\texttt{NOT} & La condición que sigue es \texttt{FALSE}
             & Mayor (se evalúa primero) \\
\texttt{AND} & \textbf{Ambas} condiciones son \texttt{TRUE}
             & Media \\
\texttt{OR}  & \textbf{Al menos una} condición es \texttt{TRUE}
             & Menor (se evalúa al final) \\
\bottomrule
\end{tabularx}
\caption{Operadores lógicos y su precedencia.}
\label{tab:operadores-logicos}
\end{table}

\begin{lstlisting}[caption={Uso de AND, OR y NOT.}]
-- AND: empleados de TI con salario mayor a 47,000
SELECT nombre, departamento, salario
FROM   empleados
WHERE  departamento = 'TI'
  AND  salario > 47000;

-- OR: empleados de Ventas O de Finanzas
SELECT nombre, departamento
FROM   empleados
WHERE  departamento = 'Ventas'
  OR   departamento = 'Finanzas';

-- NOT: todos los departamentos excepto RRHH
SELECT nombre, departamento
FROM   empleados
WHERE  NOT departamento = 'RRHH';
-- Equivalente: WHERE departamento <> 'RRHH'
\end{lstlisting}

\begin{lstlisting}[caption={Combinación de operadores y uso de paréntesis.}]
-- Gerentes de TI O analistas de Ventas con salario > 30,000
-- Los paréntesis son ESENCIALES aquí
SELECT nombre, departamento, cargo, salario
FROM   empleados
WHERE  (departamento = 'TI'     AND cargo = 'Gerente')
  OR   (departamento = 'Ventas' AND cargo = 'Analista'
                                AND salario > 30000);
\end{lstlisting}

\begin{nota}
\texttt{AND} tiene mayor precedencia que \texttt{OR}, igual que la
multiplicación tiene mayor precedencia que la suma en aritmética.
Sin paréntesis, \texttt{A OR B AND C} se interpreta como
\texttt{A OR (B AND C)}, lo que puede producir resultados
inesperados.
\textbf{Usa siempre paréntesis} cuando combines \texttt{AND} y
\texttt{OR} en la misma condición.
\end{nota}

%--------------------------------------------------------------
\subsection{\texttt{LIMIT}}
%--------------------------------------------------------------

\texttt{LIMIT} restringe el número máximo de filas que la consulta
devuelve.
Combinado con \texttt{OFFSET} (o la sintaxis abreviada de MySQL)
permite implementar \textbf{paginación}.

\begin{lstlisting}[caption={LIMIT básico.}]
-- Los 5 empleados con mayor salario
SELECT nombre, salario
FROM   empleados
ORDER BY salario DESC
LIMIT 5;
\end{lstlisting}

\begin{lstlisting}[caption={LIMIT con OFFSET para paginación.}]
-- Página 1: primeras 3 filas (OFFSET 0, implícito)
SELECT nombre, departamento
FROM   empleados
ORDER BY nombre ASC
LIMIT 3;

-- Página 2: filas 4, 5 y 6
SELECT nombre, departamento
FROM   empleados
ORDER BY nombre ASC
LIMIT 3 OFFSET 3;

-- Página 3: filas 7, 8 y 9
SELECT nombre, departamento
FROM   empleados
ORDER BY nombre ASC
LIMIT 3 OFFSET 6;
\end{lstlisting}

\begin{lstlisting}[caption={Sintaxis abreviada MySQL: LIMIT offset, cantidad.}]
-- Equivalente a LIMIT 3 OFFSET 6
SELECT nombre, departamento
FROM   empleados
ORDER BY nombre ASC
LIMIT 6, 3;    -- primero el offset, luego la cantidad
\end{lstlisting}

\begin{ejemplo}
La paginación es imprescindible en aplicaciones web.
Si una tabla tiene 500,000 registros y una página muestra 20,
la consulta para la página 3 sería:
\texttt{LIMIT 20 OFFSET 40}.
Sin \texttt{LIMIT}, traer 500,000 filas al navegador colapsaría
tanto la base de datos como la red.
\end{ejemplo}

\begin{nota}
\texttt{LIMIT} sin \texttt{ORDER BY} produce resultados
\textbf{no deterministas}: en dos ejecuciones distintas podrían
devolverse filas diferentes.
Siempre acompáñalo de \texttt{ORDER BY} para obtener páginas
consistentes.
\end{nota}



% ─── MÓDULO 5: Modificadores Parte 2 ─────────────────────
\newpage

\section{Modificadores: Parte 2}

En esta sección cubrimos los modificadores avanzados de consulta:
funciones de agregación, operadores de rango y lista, alias,
concatenación, agrupamiento, lógica condicional y manejo de nulos.
Todos son componentes esenciales para construir consultas analíticas
y reportes sobre los datos.


%--------------------------------------------------------------
\subsection{Valores \texttt{NULL}}
%--------------------------------------------------------------

\texttt{NULL} no es un valor como \texttt{0} o \texttt{''};
es la \textbf{ausencia de valor}.
Representa información desconocida o no aplicable, y su
comportamiento en operaciones es especial:
cualquier operación aritmética o comparación con \texttt{NULL}
devuelve \texttt{NULL} (ni \texttt{TRUE} ni \texttt{FALSE}).

\begin{definicion}{\texttt{NULL}}
Marcador especial que indica la ausencia de un valor en una columna.
No es equivalente a cero, a una cadena vacía ni a \texttt{FALSE}.
Para comprobar si un valor es \texttt{NULL} se usan los operadores
\texttt{IS NULL} e \texttt{IS NOT NULL}; el operador \texttt{=}
nunca devuelve \texttt{TRUE} al comparar con \texttt{NULL}.
\end{definicion}

\begin{lstlisting}[caption={Comportamiento de NULL y cómo detectarlo.}]
-- INCORRECTO: esta condición nunca devuelve filas
-- porque NULL = NULL es NULL (no TRUE)
SELECT * FROM empleados WHERE telefono = NULL;   -- ¡MAL!

-- CORRECTO: usar IS NULL
SELECT * FROM empleados WHERE telefono IS NULL;


-- Empleados que SÍ tienen teléfono registrado
SELECT * FROM empleados WHERE telefono IS NOT NULL;

-- NULL en aritmética: el resultado es siempre NULL
SELECT nombre, salario + NULL AS resultado  -- → NULL
FROM   empleados;

-- NULL en ORDER BY: MySQL los coloca al INICIO en ASC
-- y al FINAL en DESC (varía entre SGBD)
SELECT nombre, telefono
FROM   empleados
ORDER BY telefono ASC;
\end{lstlisting}

\begin{ejemplo}
Supón que agregas una columna \texttt{telefono} a la tabla
\texttt{empleados} sin indicar \texttt{NOT NULL}.
Los registros existentes quedarán con \texttt{NULL} en esa columna,
no con \texttt{0} ni con \texttt{''}.
Si después calculas \texttt{COUNT(telefono)}, esos registros
\textit{no} se contarán, porque \texttt{COUNT} ignora \texttt{NULL}s.
\end{ejemplo}

%--------------------------------------------------------------
\subsection{\texttt{MIN} y \texttt{MAX}}
%--------------------------------------------------------------

\texttt{MIN} y \texttt{MAX} son \textbf{funciones de agregación}
que devuelven el valor mínimo y máximo de una columna dentro de
un conjunto de filas.
Funcionan sobre números, fechas y cadenas de texto.

\begin{lstlisting}[caption={MIN y MAX sobre números y fechas.}]
-- Salario más bajo y más alto de toda la empresa
SELECT MIN(salario) AS salario_minimo,
       MAX(salario) AS salario_maximo
FROM   empleados;

-- MIN y MAX por departamento
SELECT departamento,
       MIN(salario) AS minimo,
       MAX(salario) AS maximo,
       MAX(salario) - MIN(salario) AS rango_salarial
FROM   empleados
GROUP BY departamento
ORDER BY rango_salarial DESC;

-- Fecha de ingreso más antigua y más reciente
SELECT MIN(fecha_ingreso) AS primer_ingreso,
       MAX(fecha_ingreso) AS ultimo_ingreso
FROM   empleados;
\end{lstlisting}

\begin{nota}
\texttt{MIN} y \texttt{MAX} \textbf{ignoran} los valores
\texttt{NULL} automáticamente.
Si todas las filas de un grupo son \texttt{NULL}, el resultado
de la función es \texttt{NULL}.
\end{nota}

%--------------------------------------------------------------
\subsection{\texttt{COUNT}}
%--------------------------------------------------------------

\texttt{COUNT} cuenta filas o valores no nulos.
Tiene tres variantes con comportamientos distintos:

\begin{lstlisting}[caption={Variantes de COUNT.}]
-- COUNT(*): cuenta TODAS las filas, incluyendo NULLs
SELECT COUNT(*) AS total_empleados
FROM   empleados;
-- Resultado: 10

-- COUNT(columna): cuenta sólo las filas donde
-- esa columna NO es NULL
SELECT COUNT(telefono) AS empleados_con_telefono
FROM   empleados;
-- Si 3 tienen NULL en telefono → resultado: 7

-- COUNT(DISTINCT columna): cuenta valores únicos no nulos
SELECT COUNT(DISTINCT departamento) AS num_departamentos
FROM   empleados;
-- Resultado: 4

-- COUNT con GROUP BY: total por grupo
SELECT departamento, COUNT(*) AS num_empleados
FROM   empleados
GROUP BY departamento
ORDER BY num_empleados DESC;
\end{lstlisting}

%--------------------------------------------------------------
\subsection{\texttt{SUM}}
%--------------------------------------------------------------

\texttt{SUM} suma todos los valores no nulos de una columna
numérica dentro del conjunto de filas evaluado.

\begin{lstlisting}[caption={SUM: suma total y por grupo.}]
-- Masa salarial total de la empresa
SELECT SUM(salario) AS masa_salarial_total
FROM   empleados;

-- Masa salarial por departamento
SELECT departamento,
       SUM(salario)  AS masa_salarial,
       COUNT(*)      AS num_empleados
FROM   empleados
GROUP BY departamento
ORDER BY masa_salarial DESC;

-- SUM condicional usando CASE (suma sólo gerentes)
SELECT SUM(CASE WHEN cargo = 'Gerente'
                THEN salario ELSE 0 END) AS total_gerentes
FROM   empleados;
\end{lstlisting}

%--------------------------------------------------------------
\subsection{\texttt{AVG}}
%--------------------------------------------------------------

\texttt{AVG} calcula el \textbf{promedio aritmético} de los valores
no nulos de una columna numérica.
Equivale a \texttt{SUM(col) / COUNT(col)}, pero ignorando
\texttt{NULL}s en ambos términos.

\begin{lstlisting}[caption={AVG: promedio general y por grupo.}]
-- Salario promedio de toda la empresa
SELECT ROUND(AVG(salario), 2) AS salario_promedio
FROM   empleados;

-- Promedio por departamento, con indicador sobre/bajo media global
SELECT departamento,
       ROUND(AVG(salario), 2) AS promedio_dpto,
       ROUND(AVG(salario) - (SELECT AVG(salario)
                             FROM empleados), 2) AS diferencia_global
FROM   empleados
GROUP BY departamento
ORDER BY promedio_dpto DESC;
\end{lstlisting}

\begin{nota}
\texttt{AVG} ignora los \texttt{NULL}s: si una columna tiene 10
filas pero 3 son \texttt{NULL}, el promedio se calcula sobre los
7 valores restantes, no dividiendo entre 10.
Si quieres tratar los \texttt{NULL}s como cero, usa
\texttt{AVG(COALESCE(columna, 0))}.
\end{nota}

%--------------------------------------------------------------
\subsection{\texttt{IN}}
%--------------------------------------------------------------

\texttt{IN} verifica si el valor de una columna pertenece a una
\textbf{lista de valores} especificada.
Es equivalente a una cadena de condiciones \texttt{OR} sobre la
misma columna, pero más legible y compacto.

\begin{lstlisting}[caption={IN con lista de valores literales.}]
-- Empleados de Ventas, TI o Finanzas
-- Equivalente: WHERE departamento = 'Ventas'
--              OR departamento = 'TI'
--              OR departamento = 'Finanzas'
SELECT nombre, departamento
FROM   empleados
WHERE  departamento IN ('Ventas', 'TI', 'Finanzas');

-- NOT IN: departamentos excluidos
SELECT nombre, departamento
FROM   empleados
WHERE  departamento NOT IN ('RRHH');

-- IN con subconsulta: departamentos que tienen gerente
SELECT nombre, departamento
FROM   empleados
WHERE  departamento IN (
    SELECT DISTINCT departamento
    FROM   empleados
    WHERE  cargo = 'Gerente'
);
\end{lstlisting}

\begin{nota}
\textbf{Precaución con \texttt{NOT IN} y \texttt{NULL}s:}
si la lista de \texttt{NOT IN} contiene un \texttt{NULL},
la condición completa devuelve \texttt{NULL} (ninguna fila),
porque SQL no puede determinar si el valor es distinto de algo
desconocido.
En subconsultas, usa \texttt{NOT EXISTS} en lugar de \texttt{NOT IN}
si los resultados pueden contener \texttt{NULL}s.
\end{nota}

%--------------------------------------------------------------
\subsection{\texttt{BETWEEN}}
%--------------------------------------------------------------

\texttt{BETWEEN} filtra filas cuyos valores se encuentran dentro
de un \textbf{rango inclusivo} (incluye los extremos).
Funciona con números, fechas y cadenas.

\begin{lstlisting}[caption={BETWEEN con números y fechas.}]
-- Empleados con salario entre 35,000 y 52,000 (inclusive)
-- Equivalente: WHERE salario >= 35000 AND salario <= 52000
SELECT nombre, salario
FROM   empleados
WHERE  salario BETWEEN 35000 AND 52000;

-- NOT BETWEEN: fuera del rango
SELECT nombre, salario
FROM   empleados
WHERE  salario NOT BETWEEN 35000 AND 52000;

-- BETWEEN con fechas
SELECT nombre, fecha_ingreso
FROM   empleados
WHERE  fecha_ingreso BETWEEN '2019-01-01' AND '2021-12-31';

-- BETWEEN con cadenas (orden alfabético)
SELECT nombre
FROM   empleados
WHERE  nombre BETWEEN 'C' AND 'M';  -- nombres de C a M
\end{lstlisting}

%--------------------------------------------------------------
\subsection{\texttt{ALIAS} (\texttt{AS})}
%--------------------------------------------------------------

Un \textbf{alias} es un nombre temporal que se asigna a una columna
o tabla dentro de una consulta.
Mejora la legibilidad del resultado y es indispensable al trabajar
con expresiones calculadas.

\begin{lstlisting}[caption={Alias de columna con AS.}]
SELECT
    nombre                          AS empleado,
    salario                         AS sueldo_base,
    salario * 12                    AS ingreso_anual,
    ROUND(salario / 30, 2)          AS sueldo_diario,
    CONCAT(nombre, ' (', cargo, ')') AS descripcion
FROM empleados;
\end{lstlisting}

\begin{lstlisting}[caption={Alias de tabla: útil con JOINs y subconsultas.}]
-- Alias de tabla: simplifica nombres largos
SELECT e.nombre, e.departamento, e.salario
FROM   empleados AS e
WHERE  e.salario > 40000;

-- La palabra AS es opcional (funciona igual sin ella)
SELECT e.nombre
FROM   empleados e
WHERE  e.cargo = 'Gerente';
\end{lstlisting}

\begin{nota}
Los alias de columna definidos en \texttt{SELECT} \textbf{no pueden
usarse} en la cláusula \texttt{WHERE} (porque \texttt{WHERE} se
evalúa antes que \texttt{SELECT}).
Sí pueden usarse en \texttt{ORDER BY} y \texttt{HAVING} en MySQL,
aunque esto no es estándar en todos los SGBD.
\end{nota}

%--------------------------------------------------------------
\subsection{\texttt{CONCAT}}
%--------------------------------------------------------------

\texttt{CONCAT} une dos o más cadenas de texto en una sola.
Es una función (no un operador) y acepta cualquier número de
argumentos.

\begin{lstlisting}[caption={CONCAT para construir cadenas compuestas.}]
-- Nombre completo con cargo entre paréntesis
SELECT CONCAT(nombre, ' — ', cargo) AS presentacion
FROM   empleados;

-- CONCAT con separador usando CONCAT_WS
-- (WithSeparator): el primer argumento es el separador
SELECT CONCAT_WS(', ', nombre, departamento, cargo) AS ficha
FROM   empleados;

-- Combinar texto literal con valores de columnas
SELECT CONCAT('El empleado ', nombre,
              ' gana $', FORMAT(salario, 2),
              ' al mes.') AS resumen
FROM   empleados
WHERE  cargo = 'Gerente';
\end{lstlisting}

\begin{nota}
En MySQL, si \textbf{cualquier} argumento de \texttt{CONCAT} es
\texttt{NULL}, el resultado completo es \texttt{NULL}.
Para evitarlo, usa \texttt{COALESCE}: \\
\texttt{CONCAT(COALESCE(nombre, 'Sin nombre'), ' — ', cargo)}.
\texttt{CONCAT\_WS} es más tolerante: ignora los argumentos
\texttt{NULL} en lugar de devolver \texttt{NULL}.
\end{nota}

%--------------------------------------------------------------
\subsection{\texttt{GROUP BY}}
%--------------------------------------------------------------

\texttt{GROUP BY} agrupa las filas que tienen valores idénticos en
las columnas especificadas, para que las funciones de agregación
(\texttt{COUNT}, \texttt{SUM}, \texttt{AVG}, etc.) operen sobre
cada grupo por separado en lugar de sobre toda la tabla.

\begin{definicion}{\texttt{GROUP BY}}
Cláusula que divide las filas del resultado de \texttt{WHERE} en
grupos, donde cada grupo comparte el mismo valor en las columnas
de agrupamiento.
Cada grupo produce exactamente \textbf{una fila} en el resultado.
\end{definicion}

\begin{lstlisting}[caption={GROUP BY con funciones de agregación.}]
-- Número de empleados y salario promedio por departamento
SELECT departamento,
       COUNT(*)           AS num_empleados,
       ROUND(AVG(salario), 2) AS promedio,
       SUM(salario)       AS masa_salarial,
       MIN(salario)       AS minimo,
       MAX(salario)       AS maximo
FROM   empleados
GROUP BY departamento
ORDER BY masa_salarial DESC;
\end{lstlisting}

\begin{lstlisting}[caption={GROUP BY con múltiples columnas.}]
-- Combinaciones únicas de departamento + cargo con conteo
SELECT departamento,
       cargo,
       COUNT(*) AS cantidad
FROM   empleados
GROUP BY departamento, cargo
ORDER BY departamento, cargo;
\end{lstlisting}

\begin{nota}
En SQL estándar (y en MySQL con \texttt{ONLY\_FULL\_GROUP\_BY}
activado), \textbf{toda columna} que aparezca en \texttt{SELECT}
y no esté dentro de una función de agregación debe estar en
\texttt{GROUP BY}.
Si no, el SGBD no sabe qué valor mostrar de entre las múltiples
filas del grupo.
\end{nota}

%--------------------------------------------------------------
\subsection{\texttt{HAVING}}
%--------------------------------------------------------------

\texttt{HAVING} filtra \textbf{grupos} (no filas individuales).
Es el equivalente de \texttt{WHERE} pero aplicado después de que
\texttt{GROUP BY} ha formado los grupos.

\begin{definicion}{\texttt{HAVING}}
Cláusula que filtra los grupos producidos por \texttt{GROUP BY},
conservando sólo aquellos cuya condición (generalmente sobre
una función de agregación) es \texttt{TRUE}.
\end{definicion}

\begin{lstlisting}[caption={HAVING: filtrar grupos por condición agregada.}]
-- Departamentos donde el salario promedio supera 45,000
SELECT departamento,
       ROUND(AVG(salario), 2) AS promedio
FROM   empleados
GROUP BY departamento
HAVING AVG(salario) > 45000;

-- Departamentos con más de 2 empleados
SELECT departamento,
       COUNT(*) AS num_empleados
FROM   empleados
GROUP BY departamento
HAVING COUNT(*) > 2;

-- Combinando WHERE (filtra filas) y HAVING (filtra grupos)
-- Departamentos con más de 1 analista
SELECT departamento,
       COUNT(*) AS num_analistas
FROM   empleados
WHERE  cargo = 'Analista'        -- primero filtra por cargo
GROUP BY departamento
HAVING COUNT(*) > 1;             -- luego filtra grupos
\end{lstlisting}

\begin{ejemplo}
La diferencia entre \texttt{WHERE} y \texttt{HAVING} puede
resumirse así:
\texttt{WHERE} descarta filas \textit{antes} de agrupar
(actúa sobre datos individuales);
\texttt{HAVING} descarta grupos \textit{después} de agrupar
(actúa sobre resultados de agregación).
Intentar usar una función de agregación en \texttt{WHERE} produce
un error de sintaxis.
\end{ejemplo}

%--------------------------------------------------------------
\subsection{\texttt{CASE}}
%--------------------------------------------------------------

\texttt{CASE} es la expresión de \textbf{lógica condicional} de SQL:
permite evaluar condiciones dentro de una consulta y devolver
valores diferentes según cuál se cumpla.
Equivale al \texttt{if/else} o al \texttt{switch} de los lenguajes
de programación.

Existen dos formas sintácticas:

\begin{lstlisting}[caption={CASE simple (compara un valor contra una lista).}]
-- Traduce el nombre del departamento al inglés
SELECT nombre,
       CASE departamento
           WHEN 'Ventas'    THEN 'Sales'
           WHEN 'TI'        THEN 'IT'
           WHEN 'RRHH'      THEN 'HR'
           WHEN 'Finanzas'  THEN 'Finance'
           ELSE 'Unknown'
       END AS department_en
FROM empleados;
\end{lstlisting}

\begin{lstlisting}[caption={CASE buscado (evalúa condiciones booleanas).}]
-- Clasifica el salario en bandas
SELECT nombre,
       salario,
       CASE
           WHEN salario >= 55000 THEN 'Banda A (Senior)'
           WHEN salario >= 45000 THEN 'Banda B (Mid)'
           WHEN salario >= 35000 THEN 'Banda C (Junior)'
           ELSE                       'Banda D (Entrada)'
       END AS banda_salarial
FROM   empleados
ORDER BY salario DESC;
\end{lstlisting}

\begin{lstlisting}[caption={CASE dentro de funciones de agregación.}]
-- Cuenta cuántos empleados hay en cada banda salarial
SELECT
    COUNT(CASE WHEN salario >= 55000 THEN 1 END) AS banda_a,
    COUNT(CASE WHEN salario BETWEEN 45000 AND 54999 THEN 1 END) AS banda_b,
    COUNT(CASE WHEN salario BETWEEN 35000 AND 44999 THEN 1 END) AS banda_c,
    COUNT(CASE WHEN salario < 35000  THEN 1 END) AS banda_d
FROM empleados;
\end{lstlisting}

%--------------------------------------------------------------
\subsection{\texttt{IFNULL} y \texttt{COALESCE}}
%--------------------------------------------------------------

Ambas funciones permiten sustituir un valor \texttt{NULL} por un
valor alternativo.
Son esenciales para producir resultados legibles cuando los datos
pueden ser incompletos.

\begin{lstlisting}[caption={IFNULL: sustituye NULL por un valor por defecto.}]
-- Si telefono es NULL, muestra 'Sin registro'
SELECT nombre,
       IFNULL(telefono, 'Sin registro') AS telefono
FROM   empleados;

-- IFNULL en cálculos: evita que NULL infecte el resultado
SELECT nombre,
       salario + IFNULL(bono, 0) AS salario_total
FROM   empleados;
\end{lstlisting}

\begin{lstlisting}[caption={COALESCE: devuelve el primer valor no NULL de una lista.}]
-- Devuelve el primer contacto disponible entre tres opciones
SELECT nombre,
       COALESCE(telefono_celular,
                telefono_fijo,
                email,
                'Sin contacto') AS contacto_preferido
FROM   empleados;
\end{lstlisting}

\begin{nota}
\texttt{IFNULL(a, b)} es equivalente a \texttt{COALESCE(a, b)},
pero \texttt{COALESCE} acepta cualquier número de argumentos y
es parte del estándar SQL.
Prefiere \texttt{COALESCE} para mayor portabilidad entre SGBD.
\end{nota}

%--------------------------------------------------------------
\subsection{Otros Modificadores}
%--------------------------------------------------------------

\subsubsection{\texttt{DISTINCT} con funciones de agregación}

\texttt{DISTINCT} puede usarse \textit{dentro} de funciones de
agregación para contar o sumar solo valores únicos:

\begin{lstlisting}[caption={DISTINCT dentro de funciones de agregación.}]
-- Cuántos cargos diferentes existen en la empresa
SELECT COUNT(DISTINCT cargo) AS cargos_distintos
FROM   empleados;

-- Suma de salarios únicos (ignora repetidos)
SELECT SUM(DISTINCT salario) AS suma_salarios_unicos
FROM   empleados;
\end{lstlisting}

\subsubsection{\texttt{ROLLUP}: subtotales y total general}

\texttt{WITH ROLLUP} agrega filas extra de subtotales a un
resultado de \texttt{GROUP BY}:

\begin{lstlisting}[caption={GROUP BY con ROLLUP para subtotales.}]
-- Masa salarial por departamento + total general al final
SELECT
    COALESCE(departamento, 'TOTAL EMPRESA') AS nivel,
    COUNT(*)      AS empleados,
    SUM(salario)  AS masa_salarial
FROM   empleados
GROUP BY departamento WITH ROLLUP;
-- La última fila tendrá NULL en departamento (el gran total)
-- COALESCE lo convierte en 'TOTAL EMPRESA'
\end{lstlisting}

\subsubsection{Expresiones aritméticas en \texttt{GROUP BY} y \texttt{HAVING}}

\begin{lstlisting}[caption={Expresiones calculadas combinadas con agrupamiento.}]
-- Año de ingreso como criterio de agrupamiento
SELECT YEAR(fecha_ingreso)   AS anio,
       COUNT(*)              AS contrataciones,
       SUM(salario)          AS masa_salarial_ese_anio
FROM   empleados
GROUP BY YEAR(fecha_ingreso)
HAVING COUNT(*) >= 2
ORDER BY anio DESC;
\end{lstlisting}

\subsubsection{\texttt{FORMAT}: formato de números en la salida}

\begin{lstlisting}[caption={FORMAT para presentación de números.}]
-- Salarios con separador de miles y 2 decimales
SELECT nombre,
       FORMAT(salario, 2)        AS salario_fmt,    -- '48,000.00'
       FORMAT(salario * 12, 0)   AS anual_fmt       -- '576,000'
FROM   empleados
ORDER BY salario DESC;
\end{lstlisting}

